\documentclass[a4paper,
fontsize=11pt,
%headings=small,
oneside,
numbers=noperiodatend,
parskip=half-,
bibliography=totoc,
final
]{scrartcl}

\usepackage[babel]{csquotes}
\usepackage{synttree}
\usepackage{graphicx}
\setkeys{Gin}{width=.4\textwidth} %default pics size

\graphicspath{{./plots/}}
\usepackage[ngerman]{babel}
\usepackage[T1]{fontenc}
%\usepackage{amsmath}
\usepackage[utf8x]{inputenc}
\usepackage [hyphens]{url}
\usepackage{booktabs} 
\usepackage[left=2.4cm,right=2.4cm,top=2.3cm,bottom=2cm,includeheadfoot]{geometry}
\usepackage[labelformat=empty]{caption} % option 'labelformat=empty]' to surpress adding "Abbildung 1:" or "Figure 1" before each caption / use parameter '\captionsetup{labelformat=empty}' instead to change this for just one caption
\usepackage{eurosym}
\usepackage{multirow}
\usepackage[ngerman]{varioref}
\setcapindent{1em}
\renewcommand{\labelitemi}{--}
\usepackage{paralist}
\usepackage{pdfpages}
\usepackage{lscape}
\usepackage{float}
\usepackage{acronym}
\usepackage{eurosym}
\usepackage{longtable,lscape}
\usepackage{mathpazo}
\usepackage[normalem]{ulem} %emphasize weiterhin kursiv
\usepackage[flushmargin,ragged]{footmisc} % left align footnote
\usepackage{ccicons} 
\setcapindent{0pt} % no indentation in captions
\usepackage{xurl} % Breaks URLs

%%%% fancy LIBREAS URL color 
\usepackage{xcolor}
\definecolor{libreas}{RGB}{112,0,0}

\usepackage{listings}

\urlstyle{same}  % don't use monospace font for urls

\usepackage[fleqn]{amsmath}

%adjust fontsize for part

\usepackage{sectsty}
\partfont{\large}

%Das BibTeX-Zeichen mit \BibTeX setzen:
\def\symbol#1{\char #1\relax}
\def\bsl{{\tt\symbol{'134}}}
\def\BibTeX{{\rm B\kern-.05em{\sc i\kern-.025em b}\kern-.08em
    T\kern-.1667em\lower.7ex\hbox{E}\kern-.125emX}}

\usepackage{fancyhdr}
\fancyhf{}
\pagestyle{fancyplain}
\fancyhead[R]{\thepage}

% make sure bookmarks are created eventough sections are not numbered!
% uncommend if sections are numbered (bookmarks created by default)
\makeatletter
\renewcommand\@seccntformat[1]{}
\makeatother

% typo setup
\clubpenalty = 10000
\widowpenalty = 10000
\displaywidowpenalty = 10000

\usepackage{hyperxmp}
\usepackage[colorlinks, linkcolor=black,citecolor=black, urlcolor=libreas,
breaklinks= true,bookmarks=true,bookmarksopen=true]{hyperref}
\usepackage{breakurl}

%meta
%meta

\fancyhead[L]{Redaktion LIBREAS\\ %author
LIBREAS. Library Ideas, 48 (2026). % journal, issue, volume.
\href{https://doi.org/10.18452/}{\color{black}https://doi.org/10.18452/}
{}} % doi 
\fancyhead[R]{\thepage} %page number
\fancyfoot[L] {\ccLogo \ccAttribution\ \href{https://creativecommons.org/licenses/by/4.0/}{\color{black}Creative Commons BY 4.0}}  %licence
\fancyfoot[R] {ISSN: 1860-7950}

\title{\LARGE{Ein Gespräch mit den LIBREAS-Gründer*innen Maxi Kindling und Ben Kaden}}% title
\author{Maxi Kindling, Ben Kaden, Linda Freyberg} % author

\setcounter{page}{1}

\hypersetup{%
      pdftitle={Ein Gespräch mit den LIBREAS-Gründer*innen Maxi Kindling und Ben Kaden},
      pdfauthor={Maxi Kindling, Ben Kaden, Linda Freyberg},
      pdfsubject={LIBREAS. Library Ideas, 48 (2026)},
      pdfkeywords={LIBREAS, Institut für Bibliotheks- und Informationswissenschaft, IBI, Open Access, Geschichte},
      pdflicenseurl={https://creativecommons.org/licenses/by/4.0/},
      pdfcopyright={CC BY 4.0 International},
      pdfcontacturl={http://libreas.eu},
      pdfurl={https://doi.org/10.18452/},
      pdfdoi={10.18452/},
      pdflang={de},
      pdfmetalang={de}
     }



\date{}
\begin{document}

\maketitle
\thispagestyle{fancyplain} 

%abstracts

%body
Geführt von Linda Freyberg am 24.03.2026

\begin{center}\rule{0.5\linewidth}{0.5pt}\end{center}

\emph{Linda Freyberg: Es geht heute um LIBREAS: Wie alles angefangen
hat, wie es weiterging und wie ihr das alles einschätzt und empfunden
habt. Wir beginnen mit der Gründung und reisen jetzt erst mal ins
Institut für Bibliothekswissenschaft in den 2000er Jahren zurück.}

\emph{Wie kam die erste Idee für LIBREAS zustande? Wer war dabei? Was
war der Kontext am Institut? Und was war eure Motivation und in welchem
Jahr war das?}

\emph{Ben Kaden:} Also, das war im Jahr 2004: Das Institut hieß damals
noch Institut für Bibliothekswissenschaft (IB). Kurz darauf wurde es
dann zum Institut für Bibliotheks- und Informationswissenschaft (IBI).
Und es hatte gerade so die Kurve bekommen, da es kurz zuvor wegen
Sparmaßnahmen beinahe geschlossen worden wäre.

So etwas traf und trifft ja hin und wieder Institute an der
Humboldt-Universität und einige sind auch verschwunden. Jemand in der
Hochschulleitung hatte damals also entschieden, dass auch das Institut
für Bibliothekswissenschaft verzichtbar wäre. Es gab Gründe dafür, auf
die müssen wir jetzt nicht eingehen. Aber es gab zugleich natürlich
reichlich Gründe dagegen. Für uns sprach zunächst einmal dagegen, dass
wir dort studierten und uns mit dem Fach identifizierten. Entsprechend
war unsere Position: Ja, wir sehen großen Sinn darin, das einzige
Institut für Bibliothekswissenschaft an einer deutschen Universität am
Leben zu erhalten. Es war, wenn man so will, halb Eigennutz, halb Stolz
aufs Fach. Es gab also ab dem Wintersemester 2003/04 Protestbewegungen,
Demonstrationen, eine Sarg-Niederlegung vor der Staatsbibliothek.

Als sich die Situation dann konsolidiert hatte und wir immer noch ein
und am Institut waren -- ich hatte ja auch eine Stelle am Institut, saß
in der Saur-Bibliothek und arbeitete für Walther Umstätter --, dachten
wir uns, wir müssten aus dem Protest etwas Nachhaltiges ableiten.
Parallel und übergeordnet setzte sich zeitgleich etwas durch, was man
Web 2.0 nannte, also die Idee, dass Leute sehr einfach und frei Inhalte
im Internet publizieren können.

Wir hatten das Thema auch auf der Agenda unseres Studienplans. Dieser
enthielt außerdem ein Seminar, in dem wir lernten, wie man Bücher macht:
\enquote{Von der Idee zum Buch} von Petra Hauke. Das war ein tolles Projekt und
es motivierte uns, selbst ein Publikationsprojekt anzugehen.

Daneben hatten wir etwas Technik und Handwerkszeug, sprich einen
Computer mit Dream\-weaver-Zugang im Institut und ein paar
HTML-Kenntnisse. Und wir hatten ein kleines Netzwerk unter anderem mit
Elisabeth Simon und Walburga Lösch. Beide führten damals einen kleinen
Verlag namens BibSpider.\footnote{\url{https://www.bibspider.de/}} Wir
kamen ins Gespräch und dann kam irgendwann der Gedanke, naja, warum
gründen wir nicht eigentlich eine eigene Zeitschrift.

Für uns als Studierende war es zu dieser Zeit sehr schwer, in etablierte
Journals reinzukommen. In unserem Fach waren es alles noch gedruckte
Zeitschriften. Open Access war in der deutschen Bibliothekswissenschaft
noch gar nicht etabliert.

Wir haben dann diese Bausteine für uns zusammengezogen und gedacht: Hier
ist der Aufbruchsgeist, wir wollen etwas für die Sichtbarkeit der
Berliner Bibliothekswissenschaft im Internet machen -- warum
eigentlich nicht gleich eine frei zugängliche elektronische Zeitschrift?
Also probierten wir es einfach aus und bastelten uns eine eigene
Plattform. Wir meinten es ernst damit, wie man im Editorial der ersten
Ausgabe immer noch nachlesen kann\footnote{Redaktion LIBREAS, "Zur
  Zeitschrift - Young Hearts Run Free. ". \emph{LIBREAS. Library Ideas},
  1 (2005). \url{https://libreas.eu/ausgabe1/001ed1.htm}}, eine
Zeitschrift zu machen, wie wir sie selbst gern lesen wollten.

Es gab auch eine weitere, ein bisschen weniger bekannte Entwicklung mit
einem verstärkenden Effekt: Walther Umstätter hatte mich als seinen
studentischen Mitarbeiter beauftragt, herauszubekommen, was eine
elektronische Zeitschrift für Bibliothekswissenschaft eigentlich
erfordern würde. Er trug sich auch mit einer Gründungsidee. Diese
verlief allerdings im Sand, weil dann, wenn man es groß aufziehen
möchte, sofort 100 Leute Einwände haben und sagen, wie es gemacht werden
sollte. Wir als Studierende waren da einerseits freier und andererseits
naiv und dreist genug, um so etwas zu übergehen. Wir wollten nicht
darauf warten, bis die Zeit insgesamt reif ist. Denn für uns war sie es
ja bereits.

\emph{Linda: Wer war bei der ursprünglichen Gruppe dabei und inwiefern
hatte das mit der Fachschaft zu tun, in der ihr ja beide auch sehr
engagiert wart und wie würdet ihr die Atmosphäre damals beschreiben?}

\emph{Maxi Kindling:} Ja, das waren aus der Fachschaftsinitiative
damals: Ben, Manuela Schulz und ich. Ich würde auch nochmal
unterstreichen, was Ben gerade gesagt hat: Im Rückblick war es eine
Zeit, in der wir erstaunlicherweise niemandem Rechenschaft ablegen
mussten. Das IB war einfach ein Ort, an dem man solche Dinge
ausprobieren konnte. Sicher gab es auch einige kritische Rückmeldungen.
Aber uns legte niemand wirklich Steine in den Weg. Ich weiß nicht, ob
man das so heute reproduzieren könnte.

Ebenfalls wichtig ist, dass wir zwar Open Access gemacht haben, aber es
nicht OA genannt haben. Wir begannen LIBREAS nicht als Open Access
Journal. Es sollte einfach frei zu lesen und auf keinen Fall mit
irgendwelchen Schranken verbunden sein. Erst 2006 entstand der Lehrstuhl
von Peter Schirmbacher. Das IBI setzte sich dann enger mit dem Thema
Open Access auseinander. Erst dann wurde uns klar, dass das, was wir da
machen, eigentlich auch einen Namen hat. Und erst zu einem späteren
Zeitpunkt beschäftigten wir uns auch mit Dingen wie Lizenzen. Anfangs
spielte das noch keine große Rolle.

Was ich vielleicht noch ergänzen möchte, ist, dass wir, Manuela und ich,
damals überhaupt keine Sorge hatten, dass wir so ein Journal nicht
gefüllt bekommen. Denn eigentlich hätten wir uns ja auch fragen können,
wo kriegen wir überhaupt Autor*innen her? Aber einerseits gab es diese
Verbindung mit dem BibSpider-Verlag, bei der wir sichergehen konnten,
dass sie sich mit darum kümmern werden, dass wir Autor*innen bekommen.
Gleichzeitig wussten wir, dass Ben Text ohne Ende produzieren kann, weil
er damals auch ganz viel im IB-Weblog schrieb, sodass irgendwie immer
klar war, an Content wird es uns nicht mangeln und auch an niemandem,
der Editorials verfasst und so weiter. Ich glaube, da trafen ganz viele
gute Umstände zusammen, so dass sich alles so umsetzen ließ, wie wir es
wollten.

Man muss auch wirklich noch einmal betonen: Das IB war von Anfang an
wirklich total aufgeschlossen. Wir sprachen natürlich auch mit diversen
Personen. Sicher waren nicht alle vorbehaltlos überzeugt oder auch nur
glücklich, was wir da ohne die richtige Zuständigkeit des Instituts als
solches machten. Aber im Grunde ließ man uns da enorm viel Freiheit und
dafür bin ich wirklich dankbar.

\emph{Ben:} Wir hatten ja auch sofort Webspace, da ich ohnehin die
Webseite des Instituts betreute. Wir konnten im Prinzip selber sogar das
Verzeichnis anlegen. Auch Michael Heinz, damaliger IT-Admin am IB,
unterstützte uns und meinte: Macht einfach, wir finden es gut, dass
jetzt so etwas passiert. Es war, glaube ich, aber auch nur so
unkompliziert möglich, weil gerade dieser Schock der Beinaheschließung
im Raum stand. So kanalisierte man das Engagement der Studierenden in
gewisser Weise. Ich kann mich an tatsächlich viele Kommentare von Leuten
erinnern, die ein bisschen sehr überlegen meinten: Ihr macht dann zwei
Ausgaben und dann habt ihr sowieso keine Lust mehr. Oder dann findet ihr
keine Inhalte. Es gab durchaus auch Skepsis.

Und wie das so ist bei Skepsis ist und bei Leuten wie uns, hat uns das
nur noch mehr motiviert. Parallel erhielten wir auch ziemlich schnell
viel Zuspruch. Die Ausgaben füllten sich und wir konnten alles Mögliche
ausprobieren. So zogen wir es, wo es sich anbot, multilingual auf und
hatten Artikel auf Russisch und Französisch in den frühen Nummern. Wir
haben Fotostrecken reingebracht. Alles, was im Web möglich war,
probierten wir zu dieser Zeit aus. Und es war ein toller Lernprozess für
uns und im Prinzip ein fortlaufendes Projektseminar, das wir für uns in
Eigenregie durchführten.

\emph{Maxi:} Vielleicht war es auch gerade deshalb erfolgreich, weil wir
keinerlei Druck hatten, dass es ein erfolgreiches Projekt werden muss,
dass so und so viele Ausgaben entstehen müssen oder dass irgendeine Form
von Reputation daraus entsteht. Es war ja für uns komplett offen, wohin
es führt und was das Ergebnis sein wird.

Wir wurden damals auch vom Lehrstuhl von Peter Schirmbacher gut
unterstützt, auch wenn dort anfangs ein bisschen Skepsis herrschte, dass
wir nur ein vergängliches und hemdsärmeliges Studierendenprojekt sind.
Wir wurden bei der technischen Ausrichtung unterstützt und ich erinnere
mich, dass wir sehr viel Support von Uwe Müller bekamen. Recht bald
konnten wir dann auch die ersten Ausgaben als Archiv auf dem edoc-Server
publizieren. PDFs auf dem edoc-Server zu haben, war für uns ein großer
Schritt. Denn so gab es eine echte Langzeitperspektive für die Dokumente
und nicht nur HTML-Seiten und selbst hochgeladene PDFs.

\emph{Ben:} Ich kann mich an den Tag erinnern. Wir mussten nämlich nach
Adlershof fahren und haben dort vor Peter Schirmbacher und seinem Team
unseren Plan vorgestellt. Das war ein bisschen komisch, weil es wie so
eine Art Projekt-Pitch war. Am Ende war es aber völlig unkompliziert.
Wir hatten da mit Peter Schirmbacher jemanden, der Dinge für uns möglich
gemacht hat.

\emph{Linda: Und was war dann im Endeffekt das Thema der ersten Ausgabe?
Und wie seid ihr an die Autor*innen herangekommen? Sind die im Institut
auf euch zugekommen womöglich?}

\emph{Ben:} Wir sprachen Leute aktiv an. Und dann hat uns Herr Umstätter
einen Beitrag geschrieben. Das erste Thema war Bibliotheksbau. Es gab in
diesen Jahren ja eine Reihe von Neubauten und für uns in der chronischen
Großbaustelle Berlin war Architektur auch so ein attraktives Thema. So
haben wir nach und nach verschiedene Sachen zusammengetragen. Elisabeth
Simon ging ebenfalls auf Tour und holte Leute dazu. Und wir selbst
schreiben auch etwas. Es gab auch einen sehr freundlichen Grafikdesigner
in unserem Umfeld. Von ihm kam das Logo-Design. \mbox{T-Shirts} gab es auch.
Zum Launch hatten wir also auch eine kleine Art Corporate Identity. Und
die Farben sind ja immer noch aktuell.

Auf dem Bibliothekartag 2005 \emph{launchten} wir dann die erste
Ausgabe. Wir waren ohnehin mit Petra Hauke in Düsseldorf, da wir den im
Seminar entstandenen Band \enquote{Bibliothekswissenschaft -- quo vadis?}
vorstellten.

Leider erinnere ich mich nicht mehr, in welcher Form der Launch geschah.
Vermutlich haben wir einfach unsere Visitenkarten verteilt. Auf der
stand mein Name und der von Elisabeth Simon, was dazu führte, dass Leute
dachten, Maxis Name sei Elisabeth.

\emph{Maxi:} Ich glaube, BibSpider als unser damaliger Partner streute
es ebenfalls in ihrem Netzwerk. In der InetBib\footnote{\url{https://www.inetbib.de/}},
die damals ein richtig großes Ding war, wurde LIBREAS auf jeden Fall
verkündet.

\emph{Ben:} Daran kann man bei der Gelegenheit auch nochmal deutlicher
erinnern. Die InetBib ist bibliothekarische Internetgeschichte. Die
Liste war zumindest von den späteren 1990ern vielleicht bis Ende der
Nullerjahre das zentrale digitale Austauschformat für die
Bibliothekswelt in Deutschland oder im DACH-Raum.

\emph{Linda: Ihr habt ihr also die erste Ausgabe auf dem Bibliothekartag
in Düsseldorf gelauncht und dann seid ihr zurückgekommen nach Berlin.}

\emph{Und wie ging es denn dann in den nächsten Jahren weiter? Die
Redaktion hat sich ja dann verändert, nehme ich an. Sind Leute
dazugekommen?}

\emph{Ben:} Das hat aber ein bisschen gedauert, wenn ich mich richtig
erinnere. Ich denke, die ersten vier, fünf Aufgaben haben wir weitgehend
im kleineren Kreis, also zu dritt, gemacht.

\emph{Maxi:} Und mit dem Verlag BibSpider. Zu diesem gehörten Walburga
Lösch, Elisabeth Simon und ganz am Anfang Julia Pagel.

\emph{Ben:} Das stimmt. Generell hat uns der Verlag sehr unterstützt. Er
besorgte uns eine ISSN, kaufte die Domain und so weiter. Er übernahm
diese Kleinstfinanzierung, die man trotzdem immer braucht. Und ein
Diktiergerät haben sie uns auch finanziert. Das habe ich noch hier.
Handys konnten zu diesem Zeitpunkt noch nicht so gut mitschneiden. Am
Ende haben wir aber nur vielleicht vier Interviews damit geführt.

\emph{Maxi:} Das ging in dem Umfang ungefähr ein Jahr, vielleicht
anderthalb Jahre so. Am Anfang produzierten wir ja auch noch mehr
Ausgaben pro Jahr. Vier Ausgaben pro Jahr wollten wir immer machen. So
ab 2006, 2007, 2008 stießen dann mehr Leute dazu, alles Studierende des
IBI.

\emph{Linda: Ja, und meint ihr, das hat dann mit dem erwähnten Kontext
mit Peter Schirmbacher und Open Access zu tun, dass da auch das
Interesse dann gewachsen ist, oder?}

\emph{Ben:} Es hatte vor allen Dingen damit zu tun, dass wir viel Zeit
in der Saur-Bibliothek verbrachten. Andere studentische Mitarbeiter wie
Boris Jacob waren auch immer präsent. Es war ein bisschen unser zweites
Wohnzimmer. Und so hat sich das dann fortgesetzt. Wir haben so immer
wieder offensiv Leute eingesponnen. Einige blieben hängen. Boris zum
Beispiel war eine unglaubliche Bereicherung, weil er auch noch mal ganz
andere Formate bedient hat. Er startete zum Beispiel eine Podcast-Reihe.
Oder LIBREAS in Second Life. LIBREAS war ja auch etwas, womit man sich
ausprobieren konnte.

Michael Seadle kam zu dieser Zeit ebenfalls ans Institut und führte die
Unterstützung weiter. Zum Beispiel bekamen wir einen sehr schönen
Redaktionslaptop.

\emph{Maxi:} Matti Stöhr und Doreen Thiede waren zu dieser Zeit auch
dabei. Sandra Lechelt machte zeitweise mit. Karsten Schuldt kam dazu.
Und irgendwann gab es auch eine Phase, wo wir LIBREAS um Leute von
außerhalb des IBI erweitern konnten, mit Heike Stadler zum Beispiel.
Aber wenn wir ehrlich sind, war es nie so, wie wir es vielleicht uns
manchmal auch gewünscht hätten, nämlich dass wir stärker aus unserer
IBI-Bubble rauskommen.

\emph{Ben:} Wir hatten dafür damals leider auch nicht so richtig die
Kontakte. Irgendwie mussten wir immer durch den Institutsrahmen
zusammenfinden. Naiko Jahn zum Beispiel kam dazu aus dem
Magister-Kolloquium. Wir haben mehr oder weniger zeitgleich unsere
Magisterarbeit geschrieben. So kam dann eins zum anderen.

Bei der BOBCATSS-Konferenz in Zadar lernten wir Bibliothekarinnen aus
Minnesota kennen, die uns für einige Jahre und einige Beiträge eine Art
amerikanische Dépendance geschaffen haben, bis sie dann in andere Berufe
wechselten und das Bibliothekswesen so ein bisschen hinter sich ließen.
Hätten wir diese Netzwerkarbeit aktiver ausgedehnt, hätten wir uns
vielleicht stärker erweitern können. Aber wir nahmen es, wie es zu uns
kam. Es war zu diesem Zeitpunkt nicht systematisch. Wer Lust zum
Mitmachen hatte, kam einfach dazu.

\emph{Maxi:} Damals hatten wir auch Zeit. Irgendwie waren wir da alle
noch in einer Phase, in der andere Sachen einfach ohne Weiteres auch
zurückgestellt werden konnten und man sich dann auf LIBREAS fokussiert
hat.

\emph{Ben:} Das kann ich bestätigen. Ich hatte für eine Weile immer
Projektstellen mit 50~\% Anteil. Teilweise am IBI, womit das dann wieder
ganz traditionell ineinander floß. Und natürlich kommt auch dazu, dass
wir deutlich jünger und familiär nicht so eingebunden waren.

\emph{Maxi:} Bis zu dem Moment, wo ein großer Teil von uns Familien
gegründet hat, habe ich das immer noch als total einfach und so nebenher
empfunden.

\emph{Linda: Aber habt ihr euch über das Profil am Anfang oder über die
Zeit dann noch stärker Gedanken gemacht, weil es ja schon eine
bibliothekswissenschaftliche Zeitschrift ist, aber auch eben diese
leichten Themen immer behandelt hat und kulturelle Phänomene aufgenommen
hat et cetera.}

\emph{Also war das so ein bewusster Prozess, dass ihr irgendwann euch so
eine Art Profil überlegt habt oder war es von Anfang an der Gedanke,
eine gewisse Offenheit zu haben und nicht etwas rein Wissenschaftliches
anzubieten? Also, wie kam die inhaltliche Ausrichtung?}

\emph{Ben:} Ich glaube, wir waren einfach an dem Punkt, dass wir das
machen wollten, worauf wir Lust hatten. So kann man es in knapper Form
fassen.

Wir haben das in die Zeitschrift eingespielt, was uns irgendwie begegnet
ist, was wir interessant fanden und das war auch bewusst so gesetzt. Das
deutsche Bibliothekswesen besaß eigentlich bereits genug
Fachzeitschriften. Unser Anspruch war nie, mit diesen Zeitschriften zu
konkurrieren. Dass wir sie oft auch nicht übermäßig mitreißend fanden,
kam dazu. Für die fachliche Weiterbildung waren sie immer relevant und
auch interessant. Aber wir wollten doch etwas anderes machen.

Man darf auch nicht vergessen, dass wir in Berlin-Mitte saßen und das
war zu dieser Zeit sehr improvisiert und aufregend. Das Tacheles lag
beispielsweise in Laufweite. Die Stadt Berlin in ihrem beständigen
Aufbruch wirkte auch auf uns zurück. Also diese Lebensweise, wenn man
Mitte 20 ist, in Berlin und mit viel Zeit an der Hand. All das, worauf
wir Lust hatten, haben wir mitunter auch in aller Unbeholfenheit
ausprobiert und was wenig überraschend auch dazu führte, dass andere
sagten, naja, das ist irgendwie so ein studentisches Ding und gar nicht
so bibliothekswissenschaftlich. Ich persönlich tat mich selbst mit dem
Label Bibliothekswissenschaft an dieser Stelle ein bisschen schwer, auch
wenn ich es sonst immer hochhielt. Aber gerade diese Zeitschrift sollte
nun aus meiner Sicht bewusst keine reine bibliothekswissenschaftliche
Zeitschrift sein und werden, sondern eher eine, die ein
Experimentierfeld ist, auf dem man alles ausprobieren kann, worauf man
gerade Lust verspürt.

\emph{Maxi:} Ja, ich würde auch sagen, dass wir durch diese
Namenserweiterung des IB, also Institut für Bibliothekswissenschaft zu
Institut für Bibliotheks- und Informationswissenschaft, eine gewisse
thematische Offenheit mitgetragen haben, die sich ein Stück weit
ebenfalls in LIBREAS widerspiegelte. Das Spektrum wurde breiter. Aus der
Redaktion hat dann der eine mal eine interessante Idee für einen Call
gehabt oder für ein Thema und ein anderes Mal die andere. Manchmal kamen
Ideen auch von außen. Wir haben uns nie irgendwie auf etwas eingegrenzt.

\emph{Ben:} Du erinnerst mich an einen wichtigen Punkt, den man hier
vielleicht auch noch einmal betrachten kann. Das ganze Fach der
Bibliothekswissenschaft war in diesen Jahren grundlegend im Umbruch. Als
ich mein Studium begann, ging es um Aussonderungsregeln,
Aufstellungssystematiken und Verbundklassifikationen. Raus ging ich dann
mit Bibliothek 2.0 und Open Access. Über den Zeitraum von fünf Jahren
hat sich das Fach in seiner Berliner Ausprägung enorm verändert.

Die traditionellen bibliotheksorganisatorischen Fragestellungen
verschwanden weitgehend und ganz viele digitale Fragestellungen kamen
hinein. Davon haben wir sehr profitiert. Parallel vollzogen sich
zahlreiche personelle Veränderungen am IBI. Neue Akteure kamen und
brachten eine neue Agenda mit. Diese Dynamik des Faches der
Bibliothekswissenschaft sowie des Bibliothekswesens, inklusive der damit
einhergehenden großen Verunsicherung des Berufsfeldes wirkte auch in
LIBREAS hinein.

Ich habe mich während meines Studiums durchweg gefragt, wie ich mich
hier eigentlich spezialisieren will? In welche Richtung will ich gehen?
Dann das, womit ich begonnen habe, war auf einmal gar kein Thema mehr.
Diese erzwungene Agilität war für LIBREAS sehr zuträglich. Zehn Jahre
zuvor hätten wir es aus verschiedenen Gründen nicht so umsetzen können.
Und zehn Jahre später vielleicht auch nicht.

Was nebenbei vielleicht auch ein Tipp für Nachwuchswissenschaftler*innen
ist: Wenn man so eine Zeitschrift und eine Visitenkarte hat, dann hat
man immer ein tolles Thema, um auf Konferenzen Menschen anzusprechen.
Dann drückt man ihnen die Karte in die Hand, erzählt, was man macht und
bittet sie um einen Beitrag. Und in einigen Fällen sagen die Leute dann
auch zu und man erweitert so sein generelles Netzwerk. Heute nutzt man
dafür vielleicht LinkedIn. Aber das gab es damals nicht. Und so ist es
häufiger passiert. Deswegen hatten wir am Anfang auch, glaube ich, eine
höhere Zahl von englischsprachigen Beiträgen, was sich jetzt
mittlerweile doch sehr aufgelöst hat, da wir auch nicht mehr so sehr in
dieser Form unterwegs und auf Konferenzen sind. Aber wir waren zum
Beispiel mehrfach bei der IFLA. Da kam es gut an, wenn man offensiv
reingeht, jung ist und sagt, wir haben hier eine eigene Zeitschrift
gegründet.

\emph{Maxi:} Wir sind auch mit den LIBREAS-T-Shirts auf den Konferenzen
unterwegs gewesen. Als ich mich nochmal rückblickend gefragt habe, was
hat denn eigentlich so Spaß gemacht an LIBREAS, dann ist das genau das,
an das ich mich sehr gerne erinnere: Man hatte immer ein Gesprächsthema,
war eigentlich völlig frei, interessante Leute anzusprechen und sie um
Beiträge zu bitten oder überhaupt einfach inhaltlich in einen Austausch
zu kommen. Und vielleicht auch manchmal ein bisschen unkonventionell.

\emph{Linda: Spannender Punkt, wie euer Branding am Anfang war und wie
es vielleicht bis heute nachwirkt. Also sowas wie jung, wild und
studentisch. Oder was wurde über euch gesagt und wart ihr dann
vielleicht gerade deswegen so cool, dass alle für LIBREAS schreiben
wollten?}

\emph{Ben:} Alle wollten ja nicht schreiben\ldots{} Ich weiß auch nicht,
ob wir wirklich \enquote{cool} waren. Wir waren vielleicht eher nerdy,
sogar ein bisschen skurril. Aber wir haben, glaube ich, damit auch etwas
gemacht, das es zu diesem Zeitpunkt sonst nicht gab. Später kamen andere
neue Zeitschriften. Andere Kolleg*innen haben ihre eigenen
Open-Access-Zeitschriften gemacht, die möglicherweise auch deutlich
seriöser und bibliothekstheoretischer oder -praktischer waren als das,
was wir produzierten. Diese Auffächerung war sicher auch notwendig und
tat der Fachkommunikation gut.

\emph{Maxi:} Ich erinnere mich aber, dass es auf jeden Fall mindestens
eine Diskussion in InetBib gab, in der sich darüber geäußert wurde, dass
LIBREAS eine Art immerwährendes studentisches Projekt ist. Diese
Wahrnehmung hielt sich sehr lange, hat uns aber nicht wirklich
geschadet.

\emph{Ben:} Den Vorgang um diese Diskussion haben wir auch gut bei uns
im Weblog dokumentiert.

Das kann man sich heute auch nicht mehr so richtig vorstellen. Wir haben
ja parallel zu der Zeitschrift auch dieses LIBREAS-Blog ziemlich
hochfrequent bestückt. Dazu kamen auch noch Veranstaltungen. Für eine
Unkonferenz gab es einen Countdown im Blog, wofür back-to-back 28
Beiträge an 28 Tagen erschienen.\footnote{\url{https://libreas.wordpress.com/tag/frei13/page/4/}}
Heute ist es fast nicht nachvollziehbar, dass man sich hinsetzt und so
eine Menge Output produziert.

\emph{Linda: Ihr wart ja dann bei vielen Kanälen eine Art Vorreiter oder
zumindest, dass ihr Formate immer schon früh ausprobiert habt. Also
Podcast wurde eben gesagt und auch das Format Blog wurde schnell sehr
intensiv genutzt.}

\emph{Und dann waren plötzlich zehn Jahre vorbei und dann 15 und dann
20. Und man stellt fest, man ist eigentlich das älteste deutschsprachige
Open-Access-Journal im bibliothekswissenschaftlichen Bereich. Wie hat
sich das dann geändert? Gab es Punkte, wo ihr gedacht habt, jetzt haben
wir ja doch etwas sehr Großes aufgebaut, was nun schon lange Bestand
hat. Habt ihr dann gedacht, jetzt müssen wir vielleicht doch mal stärker
professionalisieren, zum Beispiel eine Stelle durch Drittmittelförderung
beantragen oder so?}

\emph{Maxi:} Ja, es fehlte immer eine Person, die überhaupt Kapazitäten
hat, sich um redaktionelle Prozesse zu kümmern. Aber ich würde nicht
sagen, dass wir den Druck hatten, uns zu professionalisieren. Wenn es
doch nötig gewesen wäre, dann um die Autor*innenbetreuung und diese
ganze Kommunikation besser zu steuern. Es hat bei uns dann doch mal
länger gedauert, bis wieder jemand ins Redaktionspostfach schaute und
Mails beantwortete. An diesen Stellen wurde es dann schon irgendwann
etwas herausfordernder. Es gab diverse Vorstöße, ein Projekt auch zu
beantragen. Wir waren auch tatsächlich schon so weit, dass wir eine
halbe Antragsskizze fertig hatten. Aber irgendwie fehlten uns dann immer
auch die Kapazitäten, das zu Ende zu bringen. Daran sieht man auch, dass
irgendwann ein Punkt kommen musste, an dem ein Kapazitätsproblem spürbar
wurde.

\emph{Ben:} Aus meiner Sicht gab es noch einen anderen Punkt, der da
hinein spielte, nämlich tatsächlich die Natur der Sache selber. Hätten
wir einen DFG-Antrag gestellt und ein
\textquotesingle professionelles\textquotesingle{} Journal gemacht, dann
hätten wir uns auf einmal auch mit der Organisation von
Peer-Review-Prozessen beschäftigen müssen und mit einer anderen
inhaltlichen Aufstellung. Das hätte die Zeitschrift verändert und in
eine andere Form gebracht. Und möglicherweise waren wir noch nicht
bereit für diesen Schritt. Und wir wären eine Zeitschrift gewesen wie
viele andere.

Unser Branding der Unkonventionalität war ja, wenn schon nicht unbedingt
bewusst gewählt, so doch organisch gewachsen. Und je professioneller man
wird, je mehr Bindungen man eingeht, desto weniger Spielräume hat man in
dieser Hinsicht. Aber der Hauptpunkt war sicher: Wenn man einen Antrag
stellen möchte, dann muss man das auch durchziehen. Eine Hürde war auch,
dass wir selbst zu der Zeit auch institutionell noch nicht so verankert
waren, dass es leicht möglich gewesen wäre, selbst einen Antrag zu
stellen.

\emph{Maxi:} Es gab damals auch nicht das eine Programm zum Beispiel der
DFG, auf das wir dann gezielt da hätten einen Antrag schreiben können.
Und da wir eben kein wissenschaftliches Journal im klassischen Sinne
waren, also mit Peer Review, hätten wir es wahrscheinlich auch eher als
ein Infrastrukturvorhaben machen müssen.

Wir haben es also bis zu einem bestimmten Punkt durchaus ernsthafter
verfolgt und dann ist es auch wieder so ein bisschen im Sande verlaufen.
Außerdem gab es dann auch den Verein, den wir ja parallel gegründet
haben.

\emph{Ben:} Der Verein sollte auch die genannten Herausforderungen
auffangen oder wenigstens abfedern. Ein wenig Druck entstand auch, weil
BibSpider ab einem Punkt nicht mehr dabei war. Wir realisierten, dass
wir uns um bestimmte Dinge strukturierter kümmern müssen, zum Beispiel,
wie wir die Domain und den Webspace bezahlen. LIBREAS war, ist und
bleibt eine ehrenamtliche Angelegenheit. Aber es gibt ja trotzdem Kosten
und wir haben auch Veranstaltungen organisiert und dann stellt sich die
Frage: Was machen wir eigentlich, wenn uns jemand 200 Euro für eine
Veranstaltung spendet? Dann braucht man ein Konto. Das waren alles
Punkte, an denen sich der Bedarf für eine dahinterliegende Struktur
zeigte. Daher gründeten wir den Verein. Der läuft ja auch ganz gut und
arbeitet auch insofern in unserem kleinen Rahmen kostendeckend. Das war
das Ziel.

Und das zweite Ziel des Vereins war natürlich, Möglichkeiten zu
schaffen, doch bei Bedarf irgendwo anders andocken zu können. Das geht
mit einem Verein einfacher als mit so einer sehr losen und auch
überhaupt nicht vertraglich irgendwie eingebundenen Redaktion und
Herausgeberschaft. Es ist über einen Verein schlicht besser zu
organisieren. Da ist eine andere Bindungskraft gegeben.

\emph{Linda: Und wie war es mit der Vereinsgründung? Und vielleicht
könnt ihr ja kurz sagen, was es für Mitgliedsformen gibt, für Mitglieder
es gibt und was ihr da so gemacht habt.}

Maxi: BibSpider hatte uns organisatorisch mit angeschoben und wir hatten
auch die IBI-An\-bindung, die naturgemäß weniger wurde, je weniger wir
selbst alle am IBI tätig waren. Dennoch erschienen wir weiterhin am IBI
und das Institut hat uns unter anderem dabei unterstützt, die Marke
LIBREAS anzumelden.

Die Vereinsgründung war 2011. Bis heute gibt es drei Formen von
Mitgliedschaften. Man kann als natürliche Person Mitglied werden für
einen minimalen Beitrag von zwölf Euro. Die Idee dazu war damals, dass
wir einen Jahresbeitrag erheben wollen und pro Monat einfach einen Euro
nehmen.

Dann gibt es die Fördermitgliedschaft für 48~Euro und die
institutionelle Mitgliedschaft für 24~Euro. Dabei sind wir auch bis
heute geblieben und es hat ausgereicht. Damit konnte natürlich nie eine
Stelle oder ein Dienstleister finanziert werden oder einen finanziellen
Ausgleich bekommen. Das war auch nie das Ziel. Aber alle Unkosten und
die eine oder andere Veranstaltung und Stipendien konnten wir so
finanzieren.

\emph{Ben:} Was man aber auch noch ehrlicherweise ergänzen muss, ist,
dass sich, wenn man einen Verein zu einer Zeitschrift gründet, die im
Ehrenamt entsteht, der ehrenamtliche Aufwand verdoppelt. Man muss auf
einmal neben der Arbeit an der Zeitschrift auch Arbeit am Verein
tätigen. Und dies bedeutet nicht, dass man sich mal trifft und nett
zusammensitzt. Sondern dazu gehören der Gang zum Notar, das Erstellen
und Aushandeln von Satzungen, Kassenprüfungen und solche Dinge.

\emph{Maxi:} Wir hatten vorher tatsächlich nicht geahnt, wie viel Arbeit
der Verein selbst auch nochmal mitbringt. Keiner von uns hatte so etwas
vorher gemacht. Also galt auch da: Wir haben es einfach gemacht, ohne es
groß vorab abzuwägen.

\emph{Linda: Ja, und was für Leute sind dann Mitglieder geworden? Und
haben die auch was mitgestalten können? Oder wollten die das überhaupt
womöglich vielleicht ja auch nicht?}

\emph{Ben:} Ja, Friends and Family, kann man erst mal sagen. Wenigstens
für die Anfangszeit.

\emph{Maxi:} Auf jeden Fall, so sind wir auch mit unseren
Gründungsmitgliedern gestartet. Und der Verein hat ja das Ziel der
Förderung der bibliotheks- und informationswissenschaftlichen
Kommunikation und als Hauptkommunikationsplattformen, die unterstützt
werden, das Journal und den Blog.

Ich denke, dass die meisten Mitglieder auch deshalb Mitglieder geworden
sind, weil sie das Journal unterstützen wollen. Und weil das eben auch
kein hoher Jahresbeitrag ist, der aber trotzdem für die Sache sehr
wichtig ist und gut half. Ich glaube die Erwartung an den Verein ist
vielleicht auch, hauptsächlich das Journal zu unterstützen. Und
punktuell hatten wir auf jeden Fall Mitgestaltungsmöglichkeiten.

\emph{Ben:} Fridays For Future beziehungsweise Libraries4Future kann man
an dieser Stelle auch noch erwähnen.

\emph{Maxi:} Also ich würde sagen, wir haben schon in dem Maß, wie wir
es konnten, auch über LIBREAS hinaus Aktivitäten unterstützt. Wir haben
auch mehrfach Stipendien angeboten für Menschen, die an Konferenzen
teilnehmen wollten. Und das Echo war auch recht unterschiedlich auf die
verschiedenen Angebote.

\emph{Linda: Vieles habt ihr jetzt schon gesagt und auch, was ihr an
LIBREAS gut fandet, vor allem am Anfang. Also vielleicht wäre es doch
schön, aus eurer persönlichen Sicht zu sagen, welche Rolle LIBREAS
hatte. Es war natürlich ein Teil eurer beruflichen und auch privaten
Biografie.}

\emph{Aber könntet ihr das so zusammenfassen oder sagen, was euch das
gelehrt hat, gegeben hat, was ihr so mitgenommen habt für euer Leben,
was ihr dadurch vielleicht erst gelernt habt?}

\emph{Maxi:} Also ich würde sagen, dass LIBREAS für mich über die Jahre
immer ein fester Anker auch innerhalb der Disziplin war. Und selbst wenn
man sich beruflich oder auch aus familiären, persönlichen Gründen
vielleicht an dem einen oder anderen Punkt etwas weiter auch vom Fach
entfernt hatte, waren wir über LIBREAS immer mit dem Fach selbst und
auch mit dem Institut verbunden. So eine Konstante über so eine lange
Zeit zu haben, jetzt eben doch schon deutlich über 20 Jahre, ist sehr
wertvoll. Ich würde sagen, wir haben da auf ganz vielen Ebenen immer
auch profitiert und das ganze Netzwerk, was da so drumrum war, ist
natürlich auch bis heute großartig.

\emph{Linda: Und bei dir, Ben, was für eine Rolle hat das für dich
gespielt?}

\emph{Ben:} Ich habe unglaublich viel gelernt, vor allen Dingen auch
inhaltlich. Das Schöne mit LIBREAS und für Studierende meiner Generation
ganz spannend war, dass man, wenn man zum Beispiel auf einer Messe
herumschlenderte, zu irgendeinem Verlag gehen und sagen konnte: Wir
haben hier ein Journal, wir würden gerne das Buch rezensieren und dann
wurde es uns zugeschickt. Ich habe hier immer noch einen ganzen
Regalmeter mit Titeln, die ich zum Rezensieren bekam. Mal habe ich sie
besprochen, mal habe ich sie genossen. Damals war das fantastisch und
ich habe ganz viele Dinge gelesen, die ich sonst nicht gelesen hätte.\\
Und ich habe über diverse Dinge geschrieben und wenn man über Dinge
schreibt, lernt man auch viel über die Dinge, über die man schreibt.
Entsprechend war das für mich eine absolut prägende Phase, in der ich im
Prinzip die Themen gefunden habe, mit denen ich mich dann intensiver
beschäftigte. Über Digital Humanities habe ich zum ersten Mal in LIBREAS
geschrieben. Da hatte irgendjemand das Buch von Franco Moretti auf den
Tisch gelegt und gefragt, ob ich das nicht rezensieren will. Also habe
ich eine Rezension geschrieben und nebenbei gelernt, was Digital
Humanities eigentlich sind und welche Potentiale es da gibt. Das hat die
Agenda gesetzt. Bei Open Access war es genauso.

Diese Spezialisierung von Open Access, die viele von uns haben aus der
Redaktion, die hat sich notwendigerweise fast abgeleitet daraus, dass
wir sehr früh ein de facto Diamond Open Access Journal herausgegeben
haben. Keine Kosten für die Autor*innen, keine Kosten für die Lesenden.
Das hat mich im Prinzip dahin geführt, wo ich heute arbeite. Es war
wirklich sehr wichtig für mich, das zu machen.

\emph{Maxi:} Ja, das auf jeden Fall.

Was mir noch einfällt, ist, dass wir teilweise Themen bearbeitet haben,
die so weit weg waren von dem, was man vielleicht in der Tätigkeit als
studentische Hilfskraft gemacht hat oder dann später auf den jeweiligen
Stellen. Es gab viele sehr spezifisch bibliothekarische Themen, mit
denen ich ja sehr wenig Berührung im Alltag hatte. Aber über LIBREAS
konnte ich immer noch ein bisschen am Ball bleiben und dazu Leute
kennenlernen und, wie wir am Anfang schon gesagt haben, dann auch
einfach mal fragen: Wollt ihr nicht mal einen Beitrag für LIBREAS
darüber schreiben? Das finde sehr wertvoll.

\emph{Linda: Ihr seid ja beide als Gründungs-Herausgeber*innen und als
Redaktionsmitglieder*innen mittlerweile ausgeschieden und
voraussichtlich nur noch eine Amtsperiode im Verein aktiv. Was waren die
Gründe dafür, möchtet ihr dazu etwas sagen?}

\emph{Maxi:} Das ist wirklich verbunden mit einem lachenden und einem
weinenden Auge. Einerseits werden Kapazitäten frei und und andererseits
steckt einfach so viel Herzblut drin. Es tut schon auch weh. Aber ich
glaube, für mich ist jetzt so ein Punkt erreicht gewesen, wo es wichtig
war, auch mal Dinge loszulassen und vielleicht auch neuen Personen eine
Möglichkeit zu geben, sich zu entfalten, so wie wir es damals konnten
oder auch ganz anders. Ich glaube, es funktioniert nicht gut, wenn immer
noch sozusagen \enquote*{die Alten}, die ja auch eine bestimmte
Vorstellung davon haben, wie manche Sachen damals gemacht wurden und wie
man sie dann später weiterentwickelt hat, den Ton angeben. Somit ist es
auch gut, wenn man sich da löst und das vertrauensvoll in gute Hände
gibt. Und das können wir ja mit euch auch gerade machen.

Ben: Ich muss sagen, dass ich sehr ungern gehe. Aber ich sehe, dass es
notwendig ist. Ich habe es vorhin auch schon angedeutet. LIBREAS war ja
ein sehr junges, dynamisches Ding. Wir haben eigentlich schon frühzeitig
versucht, die nächste Generation reinzuholen. Und dann hatten wir oft
wunderbare Menschen dabei, die so für zwei, drei Jahre mitmachten und
wieder rausgegangen sind. Aber eigentlich zu selten.

Und das ist vielleicht die eine Sache, was ich auch in der Rückschau
einen Hauch bedaure, nämlich, dass wir tatsächlich häufig etwas
hermetisch gewirkt haben, dass wir es nie schafften, auch einmal ganz
andere Personen zu uns zu holen. Also auch Leute vielleicht, die wir
nicht persönlich kennen. Wir wollten es gar nicht sein, aber ich
befürchte, wir wirkten oft wie ein Insider-Projekt.

Und ich glaube, das hat leider auch unsere Außenwahrnehmung sehr
geprägt, was wiederum die Mitgestaltung durch externe Personen
erschwerte. Ich will nicht sagen, dass ich es bereue, aber ich stelle es
so in der Rückschau fest. Und ich denke, dass es notwendig ist, auch
diesen Raum jetzt wirklich und konsequent zu öffnen. Ich hoffe, es
gelingt, auch andere Leute für LIBREAS zu begeistern.

Ich gehe natürlich auch mit Wehmut. LIBREAS ist unser \enquote*{Kind}
und wir lassen es jetzt so in die Welt und haben keine Kontrolle mehr
darüber. Aber es wäre umso mehr ein Erfolg, wenn andere Leute mit der
Sache auch was anfangen können, wie wir damit etwas anfangen konnten.
Wenn sie etwas eigenes damit machen. Man merkt übrigens auch, dass man
woanders im Leben steht. Und es ist nicht mehr das Jahr 2005 oder 2006,
sondern das Jahr 2026.

\emph{Maxi:} Das \enquote*{Kind} ist volljährig, mittlerweile 21.

\emph{Ben:} Es fällt erstaunlich schwer trotz allem. Das muss man auch
sagen. Es ist gibt da auch eine emotionale Komponente. Das unterstreicht
eventuell noch einmal, dass LIBREAS nicht so ein professionelles Projekt
war wie ein professionelles Projekt es eben wäre. Sondern es ist
wirklich etwas, das mir sehr nahe geht, was extrem persönlich ist. Jetzt
muss es ohne uns weiterlaufen. Wir sind jetzt nicht mehr \enquote*{die
Eltern}, sondern die entfernten Verwandten. Vielleicht kann man das so
nehmen. Ich wünsche mir, dass ein paar Menschen hinzukommen, die
vielleicht halb so alt sind wie ich, die vielleicht genauso alt sind wie
LIBREAS. Und es neu erfinden.

%autor
\begin{center}\rule{0.5\linewidth}{0.5pt}\end{center}

\textbf{Dr.~Linda Freyberg}
(\url{https://orcid.org/0000-0002-4620-7571}) ist Mitherausgeberin und
Redakteurin der LIBREAS. Sie hat zum Thema \enquote{Ikonizität der Information.
Die Erkenntnisfunktion struktureller und gestalteter Bildlichkeit in der
digitalen Wissensorganisation} an der Leuphana Universität promoviert.
Sie ist wissenschaftliche Mitarbeiterin an der BBF \textbar{} Bibliothek
für Bildungsgeschichtliche Forschung des DIPF \textbar{}
Leibniz-Institut für Bildungsforschung und Bildungsinformation in Berlin
(\url{https://www.dipf.de/de/institut/personen/freyberg}) und hat dort
die Co-Leitung des DHELab (\url{https://bbf.dipf.de/de/dhelab}) inne.

\end{document}